% Full title: Functional Setting and the Guided-Mode Eigenvalue Problem
\section{Functional Setting and the Guided-Mode Eigenvalue Problem}
\label{sec:functional-setting}

\subsection{Geometry, material coefficient, and quasiperiodic spaces}

Let $d>0$ be the axial period and let
\[
  B:=\R\times(0,d)
\]
be the infinite period strip.  Choose two port locations $X^-<X^+$ and set
\[
  \mathcal C^0:=(X^-,X^+)\times(0,d),\qquad
  W^-:=(-\infty,X^-)\times(0,d),\qquad
  W^+:=(X^+,\infty)\times(0,d).
\]
The artificial ports are
\[
  \Gamma^-:=\{X^-\}\times(0,d),\qquad
  \Gamma^+:=\{X^+\}\times(0,d).
\]
Write $s^-= -1$ and $s^+=1$.  If $L^\pm>0$ is the transverse period of the
corresponding half-guide, its outwardly numbered cells and cell interfaces are
\begin{align*}
  \mathcal C_j^\pm
  &:=\bigl\{(x,y)\in B:(j-1)L^\pm<s^\pm(x-X^\pm)<jL^\pm\bigr\},
  &&j\in\N,\\
  \Gamma_j^\pm
  &:=\bigl\{(x,y)\in B:s^\pm(x-X^\pm)=jL^\pm\bigr\},
  &&j\in\N_0,
\end{align*}
so that $\Gamma_0^\pm=\Gamma^\pm$ and
$W^\pm=\bigcup_{j\geq1}\overline{\mathcal C_j^\pm}$ up to cell interfaces.

The center cell contains one bounded inclusion
$\Omega^0\Subset\mathcal C^0$.  A reference inclusion
$\Omega_1^\pm\Subset\mathcal C_1^\pm$ is repeated in the corresponding half-guide,
with
\[
  \Omega_j^\pm:=\Omega_1^\pm+s^\pm(j-1)L^\pm e_x,
  \qquad e_x=(1,0).
\]
All inclusion boundaries are assumed to be $C^{2}$ and separated from the cell
boundaries by a uniform positive distance.  The physical material interfaces
are
\[
  \Upsilon^0:=\partial\Omega^0,\qquad
  \Upsilon^\pm:=\bigcup_{j\geq1}\partial\Omega_j^\pm,\qquad
  \Upsilon:=\Upsilon^0\cup\Upsilon^-\cup\Upsilon^+.
\]
For the scalar TM model, let $\rho=n^2$ be the squared refractive index,
\begin{equation}
  \rho(x,y)=
  \begin{cases}
    (n^0)^2,&(x,y)\in\Omega^0,\\
    (n^\pm)^2,&(x,y)\in\Omega_j^\pm\text{ for some }j\geq1,\\
    1,&\text{otherwise},
  \end{cases}
  \qquad 0<\rho_-\leq \rho\leq \rho_+<\infty.
  \label{eq:material}
\end{equation}
More general positive piecewise-constant half-guide coefficients can be used without
changing the definitions below.

Fix a real Bloch quasimomentum $\beta\in[-\pi/d,\pi/d)$.  For any strip-like
domain $D$ having the two horizontal sides $y=0$ and $y=d$, define
\begin{equation}
  H^1_\beta(D):=
  \bigl\{u\in H^1(D):
  \left.u\right|_{y=d}
  =e^{\ii\beta d}\left.u\right|_{y=0}\bigr\},
  \label{eq:H1beta}
\end{equation}
where the boundary restrictions are understood as Sobolev traces.  This is a
closed subspace of $H^1(D)$ with the inherited norm.  The normal-derivative
quasiperiodicity is not imposed at the form level; it follows for functions in
the operator domain.

The trace spaces on a vertical port are defined by identifying the port with
$(0,d)$.  With
\[
  \beta_m:=\beta+\frac{2\pi m}{d},\qquad
  \psi_m(y):=d^{-1/2}e^{\ii\beta_my},\qquad m\in\Z,
\]
set, for $s\in\R$,
\begin{equation}
  H^s_\beta(\Gamma):=
  \left\{f=\sum_{m\in\Z}f_m\psi_m:
  \sum_{m\in\Z}\langle\beta_m\rangle^{2s}|f_m|^2<\infty\right\},
  \qquad \langle t\rangle=(1+t^2)^{1/2}.
  \label{eq:port-Sobolev}
\end{equation}
Thus $H^{-1/2}_\beta(\Gamma)$ is the anti-dual of
$H^{1/2}_\beta(\Gamma)$ under the extension of the $L^2(0,d)$ pairing.  We use
the product Cauchy space
\begin{equation}
  \mathcal H_\Gamma^\pm
  :=H^{1/2}_\beta(\Gamma^\pm)\times H^{-1/2}_\beta(\Gamma^\pm).
  \label{eq:Cauchy-space}
\end{equation}

\subsection{Guided modes and the \texorpdfstring{$\beta$}{beta}-formulation}

Fix $\beta\in[-\pi/d,\pi/d)$.  Following the formulation of
Fliss~\cite[(2.4)]{Fliss2013}, a \emph{guided mode} at wavenumber $k>0$ is a
nonzero field $u\in H^1(B)$ satisfying
\begin{equation}
  \left\{
  \begin{aligned}
    -\frac{1}{\rho}\Delta u&=k^2u&&\text{in }B,\\
    \left.u\right|_{y=d}
      &=e^{\ii\beta d}\left.u\right|_{y=0},\\
    \left.\partial_yu\right|_{y=d}
      &=e^{\ii\beta d}\left.\partial_yu\right|_{y=0}.
  \end{aligned}
  \right.
  \label{eq:guided-mode-Fliss}
\end{equation}
The condition $u\in H^1(B)$ includes square-integrability in the unbounded
$x$-direction.  Thus transverse confinement is part of the definition, while
pointwise or exponential decay will be a consequence in a spectral gap.

This paper uses the \emph{$\beta$-formulation}: $\beta$ is prescribed and the
wavenumber $k>0$ is sought.  More precisely, the problem is
\begin{equation}
  \boxed{\quad
  \text{find }k>0\text{ and }u\in D(A(\beta)),\ u\neq0,
  \quad A(\beta)u=k^2u.\quad}
  \tag{$E^\beta$}
  \label{eq:beta-formulation}
\end{equation}
\addtocounter{equation}{1}
Introduce the Sobolev space used in this formulation,
\begin{equation}
  H^1(\Delta,B):=\{u\in H^1(B):\Delta u\in L^2(B)\},
  \label{eq:H1-Delta}
\end{equation}
and set
\begin{equation}
  A(\beta):=-\frac{1}{\rho}\Delta,
  \qquad
  D(A(\beta)):=\left\{u\in H^1(\Delta,B):
  \begin{array}{l}
    \left.u\right|_{y=d}
      =e^{\ii\beta d}\left.u\right|_{y=0},\\[-1mm]
    \left.\partial_yu\right|_{y=d}
      =e^{\ii\beta d}\left.\partial_yu\right|_{y=0}
  \end{array}\right\}.
  \label{eq:Abeta}
\end{equation}
The closed-form argument of \cite[Section~3]{Fliss2013} shows that
$A(\beta)$ is self-adjoint and nonnegative in
\begin{equation}
  H:=L^2(B,\rho\,\dd x\dd y),
  \qquad \langle u,v\rangle_H:=\int_B\rho\,u\overline v\,\dd x\dd y.
  \label{eq:weighted-Hilbert-space}
\end{equation}

\begin{definition}[Fixed-$\beta$ guided modes]
\label{def:guided-space}
For fixed $\beta$ and $k>0$, the eigenspace of problem
\eqref{eq:beta-formulation} is
\begin{equation}
  \mathcal G(k,\beta)
  :=\ker(A(\beta)-k^2I).
  \label{eq:guided-space}
\end{equation}
Its nonzero members are precisely the guided modes in
\eqref{eq:guided-mode-Fliss}.  Their geometric multiplicity is
\begin{equation}
  m_{\mathrm g}(k,\beta):=\dim\mathcal G(k,\beta).
  \label{eq:geometric-multiplicity}
\end{equation}
\end{definition}

\begin{proposition}[Equivalent weak and transmission formulations]
\label{prop:weak-strong}
Suppose $\rho$ is given by \eqref{eq:material} and all interfaces are $C^2$.
A field $u$ belongs to $\mathcal G(k,\beta)$ if and only if
$u\in H^1_\beta(B)$ and
\begin{equation}
  \int_B\nabla u\cdot\nabla\overline v\,\dd x\dd y
  =k^2\int_B\rho\,u\overline v\,\dd x\dd y
  \qquad\forall v\in H^1_\beta(B).
  \label{eq:guided-weak-form}
\end{equation}
Every such field is locally $H^2$ away from the interfaces and satisfies
\begin{align}
  \Delta u+k^2\rho u&=0&&\text{in }B\setminus\Upsilon,\label{eq:strong-pde}\\
  [u]_{\Upsilon}&=0,&[\partial_\nu u]_{\Upsilon}&=0,\label{eq:TM-jumps}\\
  \left.u\right|_{y=d}
    &=e^{\ii\beta d}\left.u\right|_{y=0},&
  \left.\partial_yu\right|_{y=d}
    &=e^{\ii\beta d}\left.\partial_yu\right|_{y=0},
    \label{eq:strong-quasi}
\end{align}
where $\nu$ is either consistently chosen unit normal on each physical
interface and $[\cdot]$ denotes the trace from the $\nu$-side minus the trace
from the opposite side.  Conversely, a piecewise $H^2_{\mathrm{loc}}$ field
in $H^1(B)$ satisfying \eqref{eq:strong-pde}--\eqref{eq:strong-quasi} belongs
to $\mathcal G(k,\beta)$.
\end{proposition}
\status{background result.}
\begin{proofroadmap}
The equivalence of the operator equation and \eqref{eq:guided-weak-form}
follows from Green's identity in the weighted space $H$.  Global
$H^1$ regularity gives continuity of the Dirichlet trace across each smooth
interface; integration by parts on the two material sides shows that the
distributional interface source vanishes exactly when the two normal
derivatives agree.  The horizontal boundary terms cancel by quasiperiodicity,
and standard transmission regularity gives piecewise local $H^2$ regularity.
The converse is the same calculation in reverse.  Suitable background tools
are collected in \cite[Chapters~3--4]{McLean2000}.

Relevant internal notes are
\path{research/planning/muller-cauchy/notes/spaces.md} and
\path{research/planning/muller-cauchy/notes/trace-glue.md}.
The weakest acceptable form is distributional equivalence without asserting
global piecewise $H^2$ regularity at material junctions; such junctions are
excluded by the present smooth, separated geometry.
\end{proofroadmap}
\begin{proof}
We prove first the equivalence between the operator and weak formulations and
then recover the transmission conditions.

\emph{Step 1: the operator equation implies the weak equation.}
Let $u\in\mathcal G(k,\beta)$.  Then
$-\Delta u=k^2\rho u$ in $L^2(B)$ and both the Dirichlet and normal traces of
$u$ are $\beta$-quasiperiodic on the horizontal sides.  Take first
$v\in H^1_\beta(B)$ with bounded support in the $x$-direction.  Green's
identity on a bounded strip containing that support gives
\[
  \int_B\nabla u\cdot\nabla\overline v
  =k^2\int_B\rho u\overline v.
\]
Indeed, there are no vertical-end terms, and the terms at $y=0$ and $y=d$
cancel because the traces of $u$, $\partial_yu$, and $v$ carry the same
quasiperiodic phase.  Such test functions are dense in $H^1_\beta(B)$, so
continuity of both integrals proves \eqref{eq:guided-weak-form} for every
$v\in H^1_\beta(B)$.

\emph{Step 2: the weak equation implies the operator equation.}
Suppose that $u\in H^1_\beta(B)$ satisfies
\eqref{eq:guided-weak-form}.  Testing first with compactly supported smooth
functions shows in distributions that
\[
  -\Delta u=k^2\rho u.
\]
Since $\rho u\in L^2(B)$, this gives $u\in H^1(\Delta,B)$.  Its weak normal
traces on $y=0,d$ are therefore defined.  Applying Green's identity to a
boundedly supported $v\in H^1_\beta(B)$ and using the distributional equation
leaves only
\[
  \pair{\partial_yu|_{y=d}}{v|_{y=d}}
  -\pair{\partial_yu|_{y=0}}{v|_{y=0}}=0.
\]
Because $v|_{y=d}=e^{\ii\beta d}v|_{y=0}$ and the trace at $y=0$ can be
chosen arbitrarily in $H^{1/2}(\R)$, this identity is precisely
\[
  \partial_yu|_{y=d}
  =e^{\ii\beta d}\partial_yu|_{y=0}
  \quad\text{in }H^{-1/2}.
\]
The corresponding Dirichlet condition already follows from
$u\in H^1_\beta(B)$.  Hence $u\in D(A(\beta))$ and
$A(\beta)u=k^2u$, so $u\in\mathcal G(k,\beta)$.

\emph{Step 3: transmission conditions and local regularity.}
For a weak solution, the distributional equation found above holds with an
$L^2_{\mathrm{loc}}$ right-hand side.  Interior elliptic regularity gives
$u\in H^2_{\mathrm{loc}}$ in every material subdomain away from its boundary.
Since $u$ is one global $H^1$ function, its two Dirichlet traces on every
component of $\Upsilon$ agree.  Integrate the distributional equation by
parts on the two sides of one interface, using a test function supported in a
small interface neighbourhood.  The volume terms cancel by
\eqref{eq:strong-pde}, leaving
\[
  \pair{[\partial_\nu u]_{\Upsilon}}{\varphi}_{\Upsilon}=0
  \qquad\forall\varphi\in H^{1/2}(\Upsilon)
\]
locally on that component.  Thus both relations in
\eqref{eq:TM-jumps} hold, while \eqref{eq:strong-quasi} was established in
Steps 1--2.

\emph{Step 4: converse from the transmission formulation.}
Let now $u\in H^1(B)$ be piecewise $H^2_{\mathrm{loc}}$ and satisfy
\eqref{eq:strong-pde}--\eqref{eq:strong-quasi}.  For a
$\beta$-quasiperiodic test function with bounded $x$-support, integrate by
parts in the finitely many material pieces meeting its support.  The
Dirichlet traces define a global $H^1$ field, the interface terms cancel by
$[\partial_\nu u]_{\Upsilon}=0$, and the two horizontal terms cancel by
\eqref{eq:strong-quasi}.  This gives \eqref{eq:guided-weak-form}; density
extends it to all of $H^1_\beta(B)$.  The same calculation shows that the
piecewise Laplacians have no interface distribution, hence
$\Delta u=-k^2\rho u\in L^2(B)$.  Therefore $u\in D(A(\beta))$ and
$A(\beta)u=k^2u$, which completes the converse and the proof.
\end{proof}

\begin{remark}[TE polarization]
The draft develops only the TM model.  A scalar TE reduction has a
divergence-form principal part, for example
$-\nabla\cdot(\rho^{-1}\nabla u)=k^2u$, and continuity of the weighted normal
flux $\rho^{-1}\partial_\nu u$ replaces \eqref{eq:TM-jumps}.  This changes the form,
the conormal trace, and the integral-equation weights, so it is not treated as
a notational substitution here.
\end{remark}

\subsection{Essential spectrum, common gaps, and scan intervals}

For each $\star\in\{-,+\}$, extend the coefficient in $W^\star$
periodically in the $x$-direction to the whole strip $B$, and denote the
extension by $\rho^\star$.  The corresponding full periodic operator is
\begin{equation}
  A^\star(\beta):=-\frac{1}{\rho^\star}\Delta,\qquad
  D(A^\star(\beta)):=D(A(\beta)),
  \label{eq:background-operators}
\end{equation}
acting in $H^\star:=L^2(B,\rho^\star\,\dd x\dd y)$.  Thus
$A^\star(\beta)$ is the full periodic extension of the $\star$ half-guide,
not a half-guide operator with an artificial boundary condition at its port.

On the first cell $\mathcal C_1^\star$, introduce the outward cell coordinate
(with $X^\star=X^-$ for $\star=-$ and $X^\star=X^+$ for $\star=+$)
\[
  t^\star:=s^\star(x-X^\star)\in(0,L^\star),
  \qquad \partial_{t^\star}=s^\star\partial_x,
\]
so that $t^\star=0$ on $\Gamma_0^\star$ and $t^\star=L^\star$ on
$\Gamma_1^\star$.  For
\[
  \theta\in\mathbb B^\star
  :=\left[-\frac{\pi}{L^\star},\frac{\pi}{L^\star}\right),
\]
define the Floquet fiber in
$L^2(\mathcal C_1^\star,\rho^\star\,\dd x\dd y)$ by
\[
\begin{aligned}
  A^\star(\beta,\theta)u
  &:=-\frac{1}{\rho^\star}\Delta u,\\
  D(A^\star(\beta,\theta))
  :=\Bigl\{u\in H^1(\Delta,\mathcal C_1^\star):{}&
  \left.u\right|_{\Gamma_1^\star}
    =e^{\ii\theta L^\star}\left.u\right|_{\Gamma_0^\star},\\
  &\left.\partial_{t^\star}u\right|_{\Gamma_1^\star}
    =e^{\ii\theta L^\star}
      \left.\partial_{t^\star}u\right|_{\Gamma_0^\star},\\
  &\left.u\right|_{y=d}
    =e^{\ii\beta d}\left.u\right|_{y=0},\\
  &\left.\partial_yu\right|_{y=d}
    =e^{\ii\beta d}\left.\partial_yu\right|_{y=0}
  \Bigr\}.
\end{aligned}
\]
Here
$H^1(\Delta,\mathcal C_1^\star)
:=\{u\in H^1(\mathcal C_1^\star):\Delta u\in
L^2(\mathcal C_1^\star)\}$.  The value equalities are understood in
$H^{1/2}$ and the derivative equalities in the corresponding weak
$H^{-1/2}$ trace spaces.  Each fiber is self-adjoint and has compact
resolvent.  Denote its eigenvalues, repeated with multiplicity, by
$\lambda_n^\star(\beta,\theta)$, $n\geq1$.

Fliss~\cite[Proposition~3.1, especially equations~(3.3)--(3.4)]{Fliss2013}
gives the Floquet--Bloch band representation and continuity of the band
functions for this scalar periodic setting.  Applied separately to the two
periodic extensions, it gives
\begin{equation}
  \Sigma_\star(\beta)
  :=\sigma(A^\star(\beta))
  =\sigma_{\mathrm{ess}}(A^\star(\beta))
  =\bigcup_{n\geq1}
    \bigl\{\lambda_n^\star(\beta,\theta):
    \theta\in\mathbb B^\star\bigr\},
  \qquad \star\in\{-,+\}.
  \label{eq:background-bands}
\end{equation}
Fliss uses $\omega^2$ as the spectral parameter; in the present notation the
same parameter is $\lambda=k^2$.  The two spectra in
\eqref{eq:background-bands} must be computed independently when the periods,
cells, or material coefficients differ.

Write the nonempty bounded components of the individual positive-axis gaps as
\[
  G_m^-(\beta)
  :=\bigl(a_m^-(\beta),b_m^-(\beta)\bigr),
  \qquad
  G_n^+(\beta)
  :=\bigl(a_n^+(\beta),b_n^+(\beta)\bigr).
\]
Their common gaps are the nonempty intersections
\begin{equation}
\begin{aligned}
  G_{mn}(\beta)
  &:=G_m^-(\beta)\cap G_n^+(\beta)\\
  &=\left(
    \max\{a_m^-(\beta),a_n^+(\beta)\},
    \min\{b_m^-(\beta),b_n^+(\beta)\}
    \right),
  \label{eq:common-gap-intersection}
\end{aligned}
\end{equation}
which exist precisely when
\[
  \max\{a_m^-(\beta),a_n^+(\beta)\}
  <\min\{b_m^-(\beta),b_n^+(\beta)\}.
\]
Enumerate all such intervals as
$G_j(\beta)=(a_j(\beta),b_j(\beta))$,
$1\leq j\leq N(\beta)$, and assume below that $N(\beta)\geq1$.

For comparison, De~Nittis et al.~
\cite[Proposition~3.10 and Lemma~3.12]{DeNittisEtAl2021} prove
$\sigma_{\mathrm{ess}}(M)=\sigma(M_\ell)\cup\sigma(M_r)$
for a one-dimensional first-order Maxwell operator with two different
asymptotically periodic matrix weights.  Their Weyl essential spectrum is
characterized there by Zhislin sequences.  This is a structural precedent for
a two-ended periodic junction, but it does not itself prove the corresponding
statement for the present two-dimensional scalar second-order strip operator.
The proof in \Cref{app:essential-spectrum} adapts their localization mechanism
to the present strip and then uses the separate Floquet--Bloch descriptions
\eqref{eq:background-bands}.

\begin{theorem}[Essential spectrum for two different periodic half-guides]
\label{thm:two-ended-essential-spectrum}
\label{prop:asymmetric-essential-spectrum}
Under the geometry and coefficient assumptions of \eqref{eq:material}, the
Weyl essential spectrum of the junction operator is
\begin{equation}
  \sigma_{\mathrm{ess}}(A(\beta))
  =\Sigma_{\mathrm{ess}}(\beta)
  :=\Sigma_-(\beta)\cup\Sigma_+(\beta).
  \label{eq:projected-spectrum}
\end{equation}
In particular, no equality between the two periods or the two periodic
coefficients is required.
\end{theorem}

The detailed proof, including the local-compactness, Zhislin-sequence, and
one-sided periodic Weyl-sequence lemmas, is given in
\Cref{app:essential-spectrum}.

\begin{corollary}[Admissible scan intervals for fixed $\beta$]
\label{cor:admissible-scan-intervals}
\label{prop:admissible-scan-intervals}
If $k^2$ is an isolated $L^2$ eigenvalue of $A(\beta)$ lying outside this
essential spectrum, then
\[
  k^2\notin\Sigma_-(\beta)\cup\Sigma_+(\beta).
\]
Consequently, the candidate real-frequency scan set for ordinary discrete
gap modes is
\begin{equation}
  \mathcal K(\beta)
  :=\left\{k>0:
  k^2\notin\Sigma_-(\beta)\cup\Sigma_+(\beta)\right\}
  =\bigcup_{j=1}^{N(\beta)}
    \bigl(\sqrt{a_j(\beta)},\sqrt{b_j(\beta)}\bigr).
  \label{eq:admissible-k-scan}
\end{equation}
Equivalently, $k^2$ must lie in a common gap of the two full periodic
half-guide extensions.  This is a necessary scan condition and does not
assert that any particular common gap contains an eigenvalue.
\end{corollary}
\begin{proof}
By \Cref{thm:two-ended-essential-spectrum}, every isolated eigenvalue outside
the essential spectrum is excluded from both periodic spectra.  On the
positive spectral axis,
\[
  \bigl((0,\infty)\setminus\Sigma_-(\beta)\bigr)
  \cap
  \bigl((0,\infty)\setminus\Sigma_+(\beta)\bigr)
  =(0,\infty)\setminus
  \bigl(\Sigma_-(\beta)\cup\Sigma_+(\beta)\bigr),
\]
so the admissible spectral-parameter intervals are intersections of left and
right gaps, not their union.  For the pair $(m,n)$,
\eqref{eq:common-gap-intersection} and $\lambda=k^2$ give
\[
  k\in\left(
  \sqrt{\max\{a_m^-(\beta),a_n^+(\beta)\}},
  \sqrt{\min\{b_m^-(\beta),b_n^+(\beta)\}}
  \right),
\]
provided the strict endpoint inequality above holds.
\end{proof}

\begin{remark}[Numerical consequence]
For each fixed $\beta$, the practical scan proceeds as follows:
\begin{enumerate}[leftmargin=*]
  \item solve the left-cell Bloch problems while scanning
  $\theta\in\mathbb B^-$ and assemble $\Sigma_-(\beta)$;
  \item independently solve the right-cell Bloch problems while scanning
  $\theta\in\mathbb B^+$ and assemble $\Sigma_+(\beta)$;
  \item construct the two individual families of gaps;
  \item intersect the left and right gaps by
  \eqref{eq:common-gap-intersection};
  \item take square roots of the resulting $k^2$ intervals; and
  \item run the transmission-eigenvalue or smallest-singular-value scan only
  on those candidate $k$ intervals.
\end{enumerate}
The two periods $L^-$ and $L^+$, unit cells, and material coefficients may
all differ, so the two Bloch calculations cannot in general be reused.
The center defect can generate ordinary discrete $L^2$ localized modes only
at candidate points in these common gaps, but no existence in any candidate
interval is guaranteed.
\end{remark}

\begin{remark}[Embedded and threshold phenomena]
The scan set \eqref{eq:admissible-k-scan} targets ordinary discrete $L^2$
eigenvalues outside the essential spectrum, namely gap modes.  It does not
cover embedded eigenvalues, bound states in the continuum, threshold
eigenvalues, or band-edge resonances.  Such objects, if present, may satisfy
$k^2\in\Sigma_-(\beta)\cup\Sigma_+(\beta)$ and require a separate analysis.
Thus all isolated eigenvalues outside the essential spectrum must lie in the
common gaps; the statement is not a classification of every possible
eigenvalue.
\end{remark}

The former detailed coupled/decoupled resolvent argument, including the closed
forms, common resolvent point, Weyl theorem, and compactness route, has been
moved to
\path{archive/essential_spectrum_resolvent_argument.tex}.  It is retained
there as a technical proof attempt and is not included in this document.

\begin{proposition}[Discrete spectrum in a common gap]
\label{prop:gap-discrete-spectrum}
For every $j$ the set $\sigma(A(\beta))\cap G_j(\beta)$ consists only of isolated
eigenvalues of finite multiplicity.  Such eigenvalues can accumulate only at
$a_j(\beta)$ or $b_j(\beta)$.
\end{proposition}
\begin{proof}
By \Cref{thm:two-ended-essential-spectrum}, the interval $G_j(\beta)$ is
disjoint from the essential spectrum of the self-adjoint operator
$A(\beta)$.  Standard self-adjoint spectral theory,
used in the identical-background setting in
Fliss~\cite[Proposition~3.3]{Fliss2013}, therefore gives the claim.
\end{proof}

For the uniform estimates used in later sections, choose one common gap
$G_{j_*}(\beta)$ and a relatively compact open window
\begin{equation}
  I_\beta=(\lambda_-,\lambda_+),\qquad
  \overline{I_\beta}\subset G_{j_*}(\beta),\qquad
  \overline{I_\beta}\cap\Sigma_{\mathrm{ess}}(\beta)=\varnothing.
  \label{eq:strict-gap}
\end{equation}
The eigenproblem \eqref{eq:beta-formulation} itself is defined for every
$k>0$; the restriction $k^2\in I_\beta$ is imposed only in the later
Cauchy-relation analysis.  It guarantees simultaneously that any spectral
point is a discrete guided eigenvalue, that neither half-guide has a
propagating Bloch channel, and that the stable/unstable splitting is uniformly
separated.  To study an entire gap, one may cover it by such windows and treat
the band edges separately.

\begin{proposition}[Localization with two different half-guides]
\label{prop:asymmetric-localization}
If $k^2\in G_j(\beta)$ is an eigenvalue of $A(\beta)$, every associated
eigenfunction decays exponentially in both transverse directions.
\end{proposition}
\begin{proof}
The weighted-resolvent argument of
\cite[Proposition~3.3 and Theorem~3.5]{Fliss2013} applies separately at the two
ends.  More explicitly, the positive distances from $k^2$ to
$\sigma(A^-(\beta))$ and $\sigma(A^+(\beta))$ allow small exponential
conjugations $e^{-\alpha_-x}$ on $W^-$ and $e^{\alpha_+x}$ on $W^+$ while
keeping the two background resolvents bounded.  After a cutoff near each
port, the eigenvalue equation has a compactly supported right-hand side, so
these conjugated resolvent estimates give exponential weighted $H^1$ bounds
on the corresponding half-guides.  Joining the two estimates across the
bounded center cell proves two-sided exponential decay.  Pointwise decay then
follows from local elliptic regularity.
\end{proof}

Thus the informal condition $u(x,y)\to0$ as $x\to\pm\infty$ is a consequence
for gap eigenfunctions, not the functional definition of a guided mode.
